\documentclass[10pt,a4paper]{amsart}
\usepackage[margin=19mm]{geometry}
\usepackage{amsmath,amssymb,mathtools}
\usepackage[scr=rsfs,cal=euler]{mathalfa}
\usepackage{xfrac}
\usepackage[unicode,hidelinks,bookmarks=false]{hyperref}
\newcommand{\ie}{i.e.\ }
\newcommand{\cf}{cf.\ }
\newcommand{\resp}{respectively\ }
\newcommand{\Subsection}{\subsection}
\providecommand{\nobreakdash}{\nobreak\hbox{-}}
\setlength{\emergencystretch}{2em}
\multlinegap0pt
\usepackage{amssymb}
\usepackage{mathtools}
\usepackage[shortlabels]{enumitem}
\usepackage{xcolor}
\newtheorem{lemma}{Lemma}
\newcommand{\R}{\mathbb{R}}
\newcommand{\T}{\mathbb{T}}
\title{Rectangular Pegs on Jordan Curves of Finite $p$-Variation}
\author{Xiangfei Li, Yichen Pan}
\date{Source: arXiv:2609.04918v1 (CC BY 4.0)}
\begin{document}
\maketitle

\noindent\emph{Source and licence.} Rectangular Pegs on Jordan Curves of Finite $p$-Variation. Version arXiv:2609.04918v1 (CC BY 4.0), distributed by arXiv under the Creative Commons Attribution 4.0 International licence, http://creativecommons.org/licenses/by/4.0/, as recorded in the arXiv licence metadata for this exact version and shown on the abstract page https://arxiv.org/abs/2609.04918v1. Full licence notice: \texttt{licenses/arxiv-2609.04918-CC-BY-4.0.txt}.\par\smallskip

\noindent\textit{Editorial scope.} Counted unit: the article's lemma lem:smooth-polygon (source0.tex lines 226--233). The article's complete original proof of it is delivered with it (source0.tex lines 235--299, one \texttt{proof} environment of 65 lines). No mathematics is written, repaired or re-derived: the statement and the proof are the article's own lines, in the article's own notation, and no retained line is substituted or re-flowed.  No further source-local context is needed: every cross-reference of every retained line resolves inside the two delivered spans.  Nothing was omitted from any retained span, so the package carries no omission record.  No retained line contains a citation, so the wrapper carries no bibliography and no bibliography entry.  No retained line contains a figure environment, an image-inclusion command or drawing code, so no image is delivered and the package declares no asset dependency.  Editorial formatting only, in addition: an AMS \texttt{amsart} preamble that restates the article's own declarations and macros used by the retained lines, namely the environments \texttt{lemma} on separate counters, together with the standard packages those lines need.  The retained environments print in this document as Lemma 1 in order of appearance, whereas the article prints the same items with the numbering of its own full document.  The complete native source of the article as distributed in the arXiv e-print for this exact version is bundled unchanged as \texttt{sources/209398/source0.tex}.
\begin{lemma}\label{lem:smooth-polygon}
Let \(P\colon\T\to\R^2\) be a parametrized Jordan polygon which is affine and
nonconstant on each edge.  For every \(\eta>0\), there exists a smooth Jordan
embedding \(S\colon\T\to\R^2\) such that \(S(0)=P(0)\) and
\begin{equation}\label{eq:w11-smooth}
 \lVert S-P\rVert_{1\text{-var}}<\eta.
\end{equation}
\end{lemma}

\begin{proof}
Extend \(P\) periodically to \(\R\), let \(\rho_\delta\) be a nonnegative
periodic smooth approximate identity, and put
\[
 P_\delta=P*\rho_\delta.
\]
Since \(P\in W^{1,1}(\T;\R^2)\),
\[
 P_\delta' =P'*\rho_\delta\longrightarrow P'
 \quad\text{in }L^1(\T;\R^2).
\]
Consequently,
\begin{equation}\label{eq:w11-convolution}
 \lVert P_\delta-P\rVert_{1\text{-var}}
 =\int_{\T}|P_\delta'(t)-P'(t)|\,dt
 \longrightarrow0.
\end{equation}
Also, \(P_\delta\to P\) uniformly.

We claim that \(P_\delta\) is a smooth embedding for all sufficiently small
\(\delta\).  Away from the vertices, its derivative is the constant nonzero
velocity of an edge.  At a vertex, let \(v_-\) and \(v_+\) be the two adjacent
edge velocities.  They cannot be negative scalar multiples of each other,
since that would force the adjacent edges to overlap and contradict the
injectivity of \(P\).  Hence there is a linear functional \(\ell\) with
\[
 \ell(v_-)>0,
 \qquad
 \ell(v_+)>0.
\]
If \(\delta\) is smaller than half the least parameter length of an edge, then
near this vertex \(P_\delta'\) is a nonnegative weighted average of \(v_-\) and
\(v_+\).  Thus \(\ell(P_\delta')>0\), so \(P_\delta\) is regular and locally
injective there.  The same argument on the interiors of the finitely many
edges gives a finite collection of parameter intervals on each of which a
linear projection of \(P_\delta\) is strictly monotone, uniformly for small
\(\delta\).  Therefore there exists \(r>0\) such that
\begin{equation}\label{eq:local-inj}
 P_\delta(s)\neq P_\delta(t)
 \quad\text{whenever }0<d_{\T}(s,t)<r
\end{equation}
for all sufficiently small \(\delta\).

For pairs separated in parameter, compactness and injectivity of \(P\) give
\[
 m_r:=\min_{d_{\T}(s,t)\geq r}|P(s)-P(t)|>0.
\]
Uniform convergence implies
\(\lVert P_\delta-P\rVert_\infty<m_r/3\) for small \(\delta\), and hence
\[
 |P_\delta(s)-P_\delta(t)|
 \geq m_r-2\lVert P_\delta-P\rVert_\infty>0
\]
whenever \(d_{\T}(s,t)\geq r\).  Together with
\eqref{eq:local-inj}, this proves global injectivity.

Finally define
\[
 S(t)=P_\delta(t)-P_\delta(0)+P(0).
\]
Translation preserves regularity and injectivity, while \(S(0)=P(0)\) and
\(\lVert S-P\rVert_{1\text{-var}}\) equals the left-hand side of
\eqref{eq:w11-convolution}.  Choosing \(\delta\) sufficiently small proves
the lemma.
\end{proof}

\end{document}
